\chapter{Gastrointestinal Bleeding}

\section*{Definition}
Gastrointestinal (GI) bleeding is blood loss from anywhere in your digestive tract, from the mouth down to the anus. It is classified as upper GI bleeding when it arises in the esophagus, stomach, or the first part of the small intestine, and lower GI bleeding when it comes from the colon, rectum, or anus. Bleeding may be obvious -- vomiting blood, black tarry stool, or bright red blood from the rectum -- or hidden (occult), showing up only as unexplained anemia on a blood test.

\section*{When to See a Doctor}
Get emergency care immediately if you have any of these:
\begin{itemize}
  \item Vomiting blood or coffee-ground material
  \item Black, tarry stools or maroon stools
  \item Bright red blood from the rectum with dizziness, weakness, or fainting
  \item A rapid pulse, chest pain, shortness of breath, clammy skin, or a feeling of passing out
  \item Confusion or reduced urine output -- signs that the body may be going into shock
\end{itemize}

\section*{How It Develops}
GI bleeding begins when a blood vessel in the wall of the digestive tract is damaged, most often by an ulcer, inflammation, or injury. As blood is lost, your body tries to compensate: the heart beats faster, blood vessels in the skin and limbs narrow, and blood is redirected to vital organs. When the loss is large or rapid, blood pressure can fall and you may feel lightheaded, faint, or go into shock. Blood that stays in the stomach for a while is partially digested by acid, which is why vomited blood may look like coffee grounds and why bleeding from the stomach or upper intestine usually turns stool black and tarry (melena). Rapid bleeding in the colon passes quickly and appears bright red or maroon (hematochezia). Even slow, hidden blood loss can eventually drain your body's iron stores and produce fatigue and anemia.

\section*{Causes}
\begin{itemize}
  \item Upper GI bleeding: Peptic ulcers (the most common cause), gastritis, esophagitis, a tear at the junction of the stomach and esophagus (Mallory-Weiss tear), esophageal or stomach varices (enlarged veins from liver disease), and stomach or esophageal cancer
  \item Lower GI bleeding: Diverticular bleeding, angiodysplasia (abnormal blood vessels), colorectal polyps or cancer, inflammatory bowel disease, hemorrhoids, and anal fissures
  \item Medications: Aspirin and other nonsteroidal anti-inflammatory drugs (NSAIDs such as ibuprofen and naproxen), blood thinners (warfarin, rivaroxaban, apixaban), and antiplatelet drugs (clopidogrel)
  \item Other factors: H. pylori infection, heavy alcohol use, and advanced age raise the risk of bleeding ulcers
\end{itemize}

\section*{Related Symptoms}
\begin{itemize}
  \item Nausea and abdominal pain or cramping
  \item Weakness, fatigue, and pale skin from anemia
  \item Dizziness or lightheadedness when standing
  \item Shortness of breath and a racing heart
  \item Reduced urine output in heavier bleeding
\end{itemize}

\section*{Medical Evaluation}
\begin{itemize}
  \item You will be asked about the color and amount of blood, when it started, and whether you take aspirin, NSAIDs, or blood thinners.
  \item Vital signs are checked lying and standing to detect blood-pressure drops that suggest significant blood loss.
  \item Blood tests include a complete blood count, blood type and cross-match, clotting tests, and liver and kidney function.
  \item Endoscopy (EGD) is the main test for upper bleeding; colonoscopy examines the lower tract. Imaging and CT angiography may be used when bleeding is too fast for endoscopy.
\end{itemize}

\section*{Treatment Options}
\begin{itemize}
  \item Resuscitation first: intravenous fluids and blood transfusions to restore blood volume before any procedure.
  \item Acid-suppressing medications (intravenous proton pump inhibitors) help ulcers and gastritis stop bleeding and rebleed less.
  \item Endoscopic treatment -- cauterizing, clipping, or injecting the bleeding vessel -- often controls the bleeding without surgery.
  \item Variceal bleeding is treated with medicines that lower portal pressure (octreotide, terlipressin) and with endoscopic banding.
  \item Antibiotics are given to people with cirrhosis to prevent infection after a variceal bleed.
  \item If endoscopy cannot control the bleeding, options include embolization by an interventional radiologist or surgery.
  \item Treating the root cause -- eradicating H. pylori, healing ulcers, and reevaluating blood thinners -- prevents future bleeds.
\end{itemize}


\section*{Self-Care and Home Management}
\begin{itemize}
  \item Do not ignore blood in vomit or stool -- even a small episode warrants medical evaluation.
  \item Avoid aspirin, ibuprofen, naproxen, and alcohol, which can irritate the stomach lining and worsen bleeding.
  \item Do not stop prescribed blood thinners on your own; tell your doctor about the bleeding before changing any dose.
  \item Keep well hydrated, and follow your doctor's advice about iron supplements if anemia is present.
  \item Know that black stool can also be caused by iron pills or bismuth (Pepto-Bismol); if you have taken either, tell your doctor so they can interpret the finding correctly.
\end{itemize}

\section*{References}
\begin{thebibliography}{99}
\bibitem{gastrointestinal_bleeding:acg-ugib} Laine L, Barkun AN, Saltzman JR, et al. ACG Clinical Guideline: Upper Gastrointestinal and Ulcer Bleeding. \textit{Am J Gastroenterol}. 2021;116:899--917.
\bibitem{gastrointestinal_bleeding:acg-lgib} Strate LL, Gralnek IM. ACG Clinical Guideline: Management of Patients With Acute Lower Gastrointestinal Bleeding. \textit{Am J Gastroenterol}. 2016;111:459--474.
\bibitem{gastrointestinal_bleeding:consensus} Barkun AN, Bardou M, Kuipers EJ, et al. International consensus recommendations on the management of patients with nonvariceal upper gastrointestinal bleeding. \textit{Ann Intern Med}. 2010;152:101--113.
\bibitem{gastrointestinal_bleeding:bmj} Stanley AJ, Laine L. Management of acute upper gastrointestinal bleeding. \textit{BMJ}. 2019;364:l536.
\end{thebibliography}
